\documentclass[11pt,a4paper]{article}

\usepackage[margin=2.4cm]{geometry}
\usepackage[T1]{fontenc}
\usepackage{lmodern}
\usepackage{booktabs}
\usepackage{array}
\usepackage{listings}
\usepackage{xcolor}
\usepackage[labelsep=period]{caption}
\usepackage{parskip}

\lstset{
  basicstyle=\ttfamily\footnotesize,
  breaklines=true,
  columns=fullflexible,
  frame=single,
  rulecolor=\color{gray!45},
  backgroundcolor=\color{gray!6},
  showstringspaces=false,
  captionpos=b,
  xleftmargin=0.5em,
  framexleftmargin=0.5em
}

\newcommand{\code}[1]{\texttt{\small #1}}

\title{Incorrect gradients from a non power of two loss scale on the\\
TorchNeuron Native beta}
\author{Report for the AWS Neuron team}
\date{5 September 2026}

\begin{document}
\maketitle

\section*{Summary}

On a \code{trn2.48xlarge} running the TorchNeuron Native private beta, scaling a
scalar loss by a factor that is not a power of two before calling
\code{backward()} produces measurably wrong gradients. The forward loss is
unaffected and is bit identical to a CPU reference. Only the backward pass is
wrong, and how wrong it is depends on the scale factor. Dividing by 16 is exact,
dividing by 15 is not, and dividing by 3 sits in between. We hit this in the
textbook gradient accumulation idiom, where the loss of each micro batch is
divided by the number of micro batches before backpropagation. It silently
degraded four language model pretraining runs and cost us roughly nine hours of
wall clock time on sixteen devices before we found it. We have a workaround and
we are not blocked, but we believe the behaviour is a defect and we would like to
know whether it is expected on this build.

\section{Environment}

\begin{tabular}{@{}ll@{}}
\toprule
Instance & \code{trn2.48xlarge}, \code{i-077aedb8638c07784}, \code{us-east-2b} \\
Host operating system & Ubuntu 24.04.4 LTS \\
Driver & \code{aws-neuronx-dkms} 2.30.2.0 \\
Runtime & \code{aws-neuronx-runtime-lib} 2.34.10.0-ac18d186d \\
Collectives & \code{aws-neuronx-collectives} 2.34.10.0-74eaafac6 \\
Tools & \code{aws-neuronx-tools} 2.32.28.0-526c2b7f6 \\
Container & TorchNeuron Native private beta image, Ubuntu 24.04, Python 3.12.11 \\
Framework & \code{torch} 2.12.1, \code{torch-neuronx} 2.12.3.0.1636+5c472775 \\
Compiler & \code{neuronx-cc} 2.27.2878.0+8220f7ac \\
Logical cores & LNC equal to 2, so four logical cores per device \\
\bottomrule
\end{tabular}

\medskip

The workload is pretraining of a one billion parameter decoder only transformer
in the Llama style, with sixteen layers, model width 2048, sixteen attention
heads, feed forward width 5632, vocabulary 65536 and sequence length 2048.
Compute is bfloat16 with a float32 master copy of the weights, the optimizer is
AdamW wrapped in \code{ZeroRedundancyOptimizer}, and the model is executed with
\code{torch.compile} on the \code{neuron} backend, one compiled graph per
transformer block plus one for the embedding lookup and one for the final norm,
output projection and cross entropy. The full instance runs four such jobs at
once, each holding four devices, so each job has a world size of sixteen. The
micro batch is two sequences, gradients are accumulated over fifteen micro
batches per rank, and one optimizer step consumes 983040 tokens.

\section{What we observed in training}

Four bilingual pretraining runs were launched with a warmup, stable and decay
learning rate schedule and a target of thirty billion tokens. All four tracked
reference runs of the identical model and data, trained earlier on NVIDIA GH200
hardware, for about the first one and a half billion tokens. They then flattened.
Between two billion and two point eight billion tokens the training loss stopped
falling and in places rose slightly, while every reference run fell by about one
tenth of a nat over the same interval. An in loop bits per byte evaluation on
held out text showed the same picture, and by two point seven five billion tokens
the runs were five to thirteen percent worse than the weakest reference. A
further tell was that the gain vector of the final RMS norm grew steadily from
about 1.3 to about 1.55 in our runs, where in every healthy reference run it
shrinks over the same interval.

There were no error messages of any kind. No compiler diagnostics, no runtime
warnings, no allocator retries, no collective timeouts, no error correcting code
events in \code{sysfs}, and no loss spikes. The jobs simply learned less than
they should have.

\section{What we ruled out first}

Because the symptom looked like a data or hyperparameter problem, we checked the
obvious causes before suspecting the stack. All of the following were verified
correct and can be discounted.

\begin{itemize}
\item Learning rate schedule, optimizer betas, epsilon, weight decay and gradient
      clipping value, all compared key by key against the reference configuration.
\item Data mixing, per rank data disjointness, and the token budget accounting.
\item Checkpoint save and resume, including the per rank optimizer state shards.
\item The evaluation itself. We rescored a saved checkpoint on the host in float32
      with an independent implementation and reproduced the on device evaluation
      numbers, which showed the evaluation was reporting the truth about the weights.
\item Hardware health. All sixteen devices report zero corrected and zero
      uncorrected memory and static memory error counts, and no unhealthy components.
\item Forward numerics. On a fixed micro batch, the loss computed through the
      compiled training path agrees with a float32 CPU reference to four decimal
      places, and repeating the same micro batch returns a bit identical value.
\item Single micro batch gradients. With no loss scaling at all, gradients from the
      compiled path agree with the float32 CPU reference to a worst case cosine
      similarity of 0.998, which is tighter than a bfloat16 autocast CPU run
      achieves against the same reference.
\item The optimizer. Gradient clipping on device returns the same total norm we
      compute on the host, and one AdamW step taken on device agrees with one taken
      on the host from the same gradients to a cosine similarity of 1.000.
\end{itemize}

\section{The defect}

The remaining difference between the healthy single micro batch case and the
training loop was the loss scaling. The training loop calls
\code{(loss / grad\_accum).backward()} with \code{grad\_accum} equal to fifteen,
which is the standard way to average micro batches. We therefore measured
gradient accuracy as a function of that scale.

The experiment holds everything else fixed. One real micro batch of two sequences
of 2048 tokens, one saved checkpoint, one forward pass through the compiled
training path, then \code{backward()} with a scalar scale $s$ applied to the loss.
The resulting float32 master gradients are divided by $s$ and compared, parameter
by parameter, with a float32 CPU reference gradient of the same batch and the same
weights. Since dividing the loss by $s$ and multiplying the gradient by $1/s$ are
exactly inverse operations in exact arithmetic, every row of the table below
should be identical.

\begin{table}[h]
\centering
\caption{Gradient accuracy against a float32 CPU reference as a function of the
scalar loss scale. Cosine similarity is the worst value over all parameter
families. The norm ratio column is the range, over the sixteen layers, of the
ratio of the RMS norm gain gradient norm to the reference, where a correct result
is one.}
\begin{tabular}{@{}lccl@{}}
\toprule
Loss scale before \code{backward()} & Worst cosine & RMS norm gain norm ratio & Verdict \\
\midrule
1 (no scaling)   & 0.9983 & 0.96 to 1.03 & correct \\
1/2              & 0.9983 & 0.96 to 1.03 & correct \\
1/4              & 0.9983 & 0.96 to 1.03 & correct \\
1/16             & 0.9983 & 0.96 to 1.03 & correct \\
1/3              & 0.9639 & 0.76 to 1.14 & wrong \\
\textbf{1/15}    & \textbf{0.9041} & \textbf{0.53 to 1.27} & \textbf{wrong} \\
\bottomrule
\end{tabular}
\end{table}

The error is not uniform across the network. At a scale of one fifteenth the
worst affected parameters are the small float32 gain vectors of the RMS norm
layers, whose gradient norms are wrong by up to a factor of two in either
direction, and the attention value projections, whose worst cosine similarity is
0.913. The matrices are wrong by a few percent. The output projection is barely
affected. Because the error is direction dependent rather than a uniform rescale,
gradient clipping and the Adam second moment do not absorb it, and it accumulates
over training.

Three further observations narrow the cause.

\begin{itemize}
\item The forward loss is unchanged and bit identical across all scales. Only the
      backward pass differs.
\item It makes no difference how the scale is applied. Writing \code{loss / 15},
      writing \code{loss * (1.0/15)}, and passing the scale as the
      \code{gradient} argument of \code{backward()} all produce the same wrong
      numbers, to every digit we printed.
\item Repeating the same micro batch is bit deterministic, so this is not noise
      or a race. It is a reproducible function of the scale.
\end{itemize}

The pattern across scales is consistent with a value that is exactly
representable being handled correctly and a value that is not being handled
incorrectly. One half, one quarter and one sixteenth are exact in binary
floating point. One third and one fifteenth are not.

We also confirmed the practical consequence for gradient accumulation. Running the
same micro batch three times with a scale of one third gives a result whose worst
family cosine similarity against three times the correctly scaled single batch
gradient is 0.868, and the RMS norm gain gradient norms are off by a factor of
0.61 to 1.26. Accumulating two different micro batches and swapping their order
is exact, so the accumulation machinery itself is sound. The scale alone carries
the error.

\section{Minimal reproduction}

The listing below is the shape of the check we ran. It needs one device and about
a minute once the graph is compiled. Our full diagnostic scripts, which also cover
accumulation order, the optimizer step and the compile before load case, are
available on request.

\begin{lstlisting}[language=Python]
import torch, torch.nn.functional as F, torch_neuronx, torch.distributed as dist

dist.init_process_group(backend="neuron")
dev = torch.device(f"neuron:{torch_neuronx.current_device()}")

model  = build_model().to(dev)          # bf16 shadow weights, fp32 masters
blocks = [b.compile(backend="neuron", dynamic=False) for b in model.layers]

x, y = fixed_micro_batch()              # 2 sequences of 2048 tokens, on device

def grads_at(scale):
    zero_all_grads(model)
    with torch.autocast("neuron", dtype=torch.bfloat16):
        loss = forward_loss(x, y)       # scalar cross entropy
    (loss * scale).backward()
    return {k: p.grad.double().cpu() / scale for k, p in model.named_parameters()}

ref = grads_at(1.0)                     # also matches a float32 CPU reference
for s in (0.5, 0.25, 1/16, 1/3, 1/15):
    g = grads_at(s)
    worst = min(cosine(g[k], ref[k]) for k in ref)
    print(s, worst)                     # 1.000 for the powers of two, 0.90 for 1/15
\end{lstlisting}

\section{Workaround}

We removed the loss scaling from the backward pass. Each micro batch is now
backpropagated unscaled, and the division by the accumulation count is applied to
the float32 master gradients afterwards, before the all reduce across ranks.

\begin{lstlisting}[language=Python]
for x, y in micro_batches:
    with torch.autocast("neuron", dtype=torch.bfloat16):
        loss = forward_loss(x, y)
    loss.backward()                     # unscaled
for p in master.parameters():
    p.grad.div_(grad_accum)             # fp32 division on the accumulated gradient
    dist.all_reduce(p.grad, op=dist.ReduceOp.AVG)
\end{lstlisting}

With this change, accumulating fifteen copies of one micro batch and dividing by
fifteen reproduces the single micro batch gradient to a cosine similarity of
1.00000, and accumulating three different micro batches and dividing by three
reproduces their mean to the same precision. The restarted runs now track the
reference runs correctly at every budget we have reached.

\section{Questions for the Neuron team}

\begin{enumerate}
\item Is a non power of two scalar multiplier on the loss, or equivalently on the
      \code{gradient} argument of \code{backward()}, known to be unsafe on this
      build of the native backend. If so, we would suggest documenting it, since
      dividing the loss by the number of micro batches is the standard gradient
      accumulation idiom and appears in most training examples.
\item Is the error in the compiler when it lowers the seeded gradient, or in the
      runtime. We are happy to run any instrumented build or to share the compiled
      artifacts for the two cases.
\item Would a diagnostic help here. A warning when a backward is seeded with a
      value that cannot be represented exactly would have saved us the entire
      investigation, because nothing in the logs, the counters or the loss curve
      pointed at the backward pass.
\item Is our workaround the recommended one, or is there a preferred way to
      average micro batches on this stack.
\end{enumerate}

\section{Impact}

The defect is silent. It does not raise, it does not warn, and it leaves a loss
curve that looks plausible for the first half of the affected range. We lost about
three billion tokens of training on each of four concurrent jobs, roughly nine
hours of wall clock time on all sixteen devices, and about a day of engineering
time. Anyone porting a standard PyTorch training loop to this backend will write
the same line of code that we wrote.

\end{document}
